% Options for packages loaded elsewhere
\PassOptionsToPackage{unicode}{hyperref}
\PassOptionsToPackage{hyphens}{url}
%
\documentclass[
  12pt,
]{article}
\usepackage{lmodern}
\usepackage{amssymb,amsmath}
\usepackage{ifxetex,ifluatex}
\ifnum 0\ifxetex 1\fi\ifluatex 1\fi=0 % if pdftex
  \usepackage[T1]{fontenc}
  \usepackage[utf8]{inputenc}
  \usepackage{textcomp} % provide euro and other symbols
\else % if luatex or xetex
  \usepackage{unicode-math}
  \defaultfontfeatures{Scale=MatchLowercase}
  \defaultfontfeatures[\rmfamily]{Ligatures=TeX,Scale=1}
\fi
% Use upquote if available, for straight quotes in verbatim environments
\IfFileExists{upquote.sty}{\usepackage{upquote}}{}
\IfFileExists{microtype.sty}{% use microtype if available
  \usepackage[]{microtype}
  \UseMicrotypeSet[protrusion]{basicmath} % disable protrusion for tt fonts
}{}
\makeatletter
\@ifundefined{KOMAClassName}{% if non-KOMA class
  \IfFileExists{parskip.sty}{%
    \usepackage{parskip}
  }{% else
    \setlength{\parindent}{0pt}
    \setlength{\parskip}{6pt plus 2pt minus 1pt}}
}{% if KOMA class
  \KOMAoptions{parskip=half}}
\makeatother
\usepackage{xcolor}
\IfFileExists{xurl.sty}{\usepackage{xurl}}{} % add URL line breaks if available
\IfFileExists{bookmark.sty}{\usepackage{bookmark}}{\usepackage{hyperref}}
\hypersetup{
  pdftitle={Arrays from R to C and Back},
  pdfauthor={Jan de Leeuw},
  hidelinks,
  pdfcreator={LaTeX via pandoc}}
\urlstyle{same} % disable monospaced font for URLs
\usepackage[margin=1in]{geometry}
\usepackage{color}
\usepackage{fancyvrb}
\newcommand{\VerbBar}{|}
\newcommand{\VERB}{\Verb[commandchars=\\\{\}]}
\DefineVerbatimEnvironment{Highlighting}{Verbatim}{commandchars=\\\{\}}
% Add ',fontsize=\small' for more characters per line
\usepackage{framed}
\definecolor{shadecolor}{RGB}{248,248,248}
\newenvironment{Shaded}{\begin{snugshade}}{\end{snugshade}}
\newcommand{\AlertTok}[1]{\textcolor[rgb]{0.94,0.16,0.16}{#1}}
\newcommand{\AnnotationTok}[1]{\textcolor[rgb]{0.56,0.35,0.01}{\textbf{\textit{#1}}}}
\newcommand{\AttributeTok}[1]{\textcolor[rgb]{0.77,0.63,0.00}{#1}}
\newcommand{\BaseNTok}[1]{\textcolor[rgb]{0.00,0.00,0.81}{#1}}
\newcommand{\BuiltInTok}[1]{#1}
\newcommand{\CharTok}[1]{\textcolor[rgb]{0.31,0.60,0.02}{#1}}
\newcommand{\CommentTok}[1]{\textcolor[rgb]{0.56,0.35,0.01}{\textit{#1}}}
\newcommand{\CommentVarTok}[1]{\textcolor[rgb]{0.56,0.35,0.01}{\textbf{\textit{#1}}}}
\newcommand{\ConstantTok}[1]{\textcolor[rgb]{0.00,0.00,0.00}{#1}}
\newcommand{\ControlFlowTok}[1]{\textcolor[rgb]{0.13,0.29,0.53}{\textbf{#1}}}
\newcommand{\DataTypeTok}[1]{\textcolor[rgb]{0.13,0.29,0.53}{#1}}
\newcommand{\DecValTok}[1]{\textcolor[rgb]{0.00,0.00,0.81}{#1}}
\newcommand{\DocumentationTok}[1]{\textcolor[rgb]{0.56,0.35,0.01}{\textbf{\textit{#1}}}}
\newcommand{\ErrorTok}[1]{\textcolor[rgb]{0.64,0.00,0.00}{\textbf{#1}}}
\newcommand{\ExtensionTok}[1]{#1}
\newcommand{\FloatTok}[1]{\textcolor[rgb]{0.00,0.00,0.81}{#1}}
\newcommand{\FunctionTok}[1]{\textcolor[rgb]{0.00,0.00,0.00}{#1}}
\newcommand{\ImportTok}[1]{#1}
\newcommand{\InformationTok}[1]{\textcolor[rgb]{0.56,0.35,0.01}{\textbf{\textit{#1}}}}
\newcommand{\KeywordTok}[1]{\textcolor[rgb]{0.13,0.29,0.53}{\textbf{#1}}}
\newcommand{\NormalTok}[1]{#1}
\newcommand{\OperatorTok}[1]{\textcolor[rgb]{0.81,0.36,0.00}{\textbf{#1}}}
\newcommand{\OtherTok}[1]{\textcolor[rgb]{0.56,0.35,0.01}{#1}}
\newcommand{\PreprocessorTok}[1]{\textcolor[rgb]{0.56,0.35,0.01}{\textit{#1}}}
\newcommand{\RegionMarkerTok}[1]{#1}
\newcommand{\SpecialCharTok}[1]{\textcolor[rgb]{0.00,0.00,0.00}{#1}}
\newcommand{\SpecialStringTok}[1]{\textcolor[rgb]{0.31,0.60,0.02}{#1}}
\newcommand{\StringTok}[1]{\textcolor[rgb]{0.31,0.60,0.02}{#1}}
\newcommand{\VariableTok}[1]{\textcolor[rgb]{0.00,0.00,0.00}{#1}}
\newcommand{\VerbatimStringTok}[1]{\textcolor[rgb]{0.31,0.60,0.02}{#1}}
\newcommand{\WarningTok}[1]{\textcolor[rgb]{0.56,0.35,0.01}{\textbf{\textit{#1}}}}
\usepackage{graphicx}
\makeatletter
\def\maxwidth{\ifdim\Gin@nat@width>\linewidth\linewidth\else\Gin@nat@width\fi}
\def\maxheight{\ifdim\Gin@nat@height>\textheight\textheight\else\Gin@nat@height\fi}
\makeatother
% Scale images if necessary, so that they will not overflow the page
% margins by default, and it is still possible to overwrite the defaults
% using explicit options in \includegraphics[width, height, ...]{}
\setkeys{Gin}{width=\maxwidth,height=\maxheight,keepaspectratio}
% Set default figure placement to htbp
\makeatletter
\def\fps@figure{htbp}
\makeatother
\setlength{\emergencystretch}{3em} % prevent overfull lines
\providecommand{\tightlist}{%
  \setlength{\itemsep}{0pt}\setlength{\parskip}{0pt}}
\setcounter{secnumdepth}{5}
\newlength{\cslhangindent}
\setlength{\cslhangindent}{1.5em}
\newenvironment{cslreferences}%
  {\setlength{\parindent}{0pt}%
  \everypar{\setlength{\hangindent}{\cslhangindent}}\ignorespaces}%
  {\par}

\title{Arrays from R to C and Back}
\author{Jan de Leeuw}
\date{First created on May 21, 2020. Last update on June 02, 2020}

\begin{document}
\maketitle
\begin{abstract}
For multidimensional arrays and symmetric arrays in compact storage the
mapping of R array indices to C storage locations and its inverse can be
complicated. This short note discusses and implements a systematic way
of handling this mapping for vectors, general matrices, symmetric
matrices in compact storage, general multidimensional arrays, and
super-symmetric multidimensional arrays.
\end{abstract}

{
\setcounter{tocdepth}{4}
\tableofcontents
}
\begin{center}\rule{0.5\linewidth}{0.5pt}\end{center}

Note: This is a working paper which will be expanded/updated frequently.
All suggestions for improvement are welcome. The directory
\href{http://deleeuwpdx.net/pubfolders/array}{deleeuwpdx.net/pubfolders/array}
has a pdf version, the bib file, the complete Rmd file with the code
chunks, and the R and C source code.

\hypertarget{introduction}{%
\section{Introduction}\label{introduction}}

Arrays in R are vectors with a dim attribute. Arrays in C are blocks of
consecutive memory, and references to arrays \emph{usually} decay to
pointers to the initial element of the block. Arrays of pointers can be
used to create, or maybe we should say simulate, multidimensional
arrays.

There are two main differences between arrays in R and C. The first is
that array indexing in R starts with 1, and in C it starts with 0. Thus
in an R array x the first element is \(x[1]\), while in C it is \(x[0]\)
or, using pointers, \(*x\) or \(*(x + 0)\). The second difference is
that in R column-major ordering is used to store the content of arrays
of rank two or higher, while C uses row-major storage. Thus if x is the
matrix

\begin{verbatim}
##      [,1] [,2]
## [1,]   1    3 
## [2,]   2    4
\end{verbatim}

then for the third element in R storage we have \(x[3]=x[1,2]=3\), while
in C for the third element we have \(*(x + 2)=x[2][1]=2\).

If we pass arrays from R to C (or back) then this can become annoying,
especially for arrays of high rank. Even more complications arise if we
use compact storage of symmetric arrays, because the mapping of, say, R
array indices to C storage locations and its inverse are bound to be
more complicated. This short note discusses one systematic way to easily
handle that mapping in some of the more important cases.

\hypertarget{arrays-in-c}{%
\section{Arrays in C}\label{arrays-in-c}}

In array.h we define an array as a C structure with the following
elements:

\begin{itemize}
\tightlist
\item
  content - pointer to data of type CTYPE (defaults to double);
\item
  length - pointer to the number of elements in content;
\item
  rank - pointer to number of dimensions of array;
\item
  shape - pointer to vector with number of levels in each dimension;
\item
  encode - pointer to a mapping of C memory locations to R array
  indices;
\item
  decode - pointer to a mapping of R array indices to C memory
  locations.
\end{itemize}

CTYPE is a macro which you can change in array.h from double to other
types. encode and decode are of type ENCODE and DECODE, which are
typedefs for function pointers.

\begin{Shaded}
\begin{Highlighting}[]
\KeywordTok{typedef} \DataTypeTok{void}\NormalTok{ (*ENCODE)(}
    \DataTypeTok{const} \DataTypeTok{int}\NormalTok{ *, }\DataTypeTok{int}\NormalTok{ *, }\DataTypeTok{const} \DataTypeTok{int}\NormalTok{ *,}
    \DataTypeTok{const} \DataTypeTok{int}\NormalTok{ *); }\CommentTok{// encode gives R array indices from C memory location}
\KeywordTok{typedef} \DataTypeTok{void}\NormalTok{ (*DECODE)(}
    \DataTypeTok{const} \DataTypeTok{int}\NormalTok{ *, }\DataTypeTok{int}\NormalTok{ *, }\DataTypeTok{const} \DataTypeTok{int}\NormalTok{ *,}
    \DataTypeTok{const} \DataTypeTok{int}\NormalTok{ *); }\CommentTok{// decode gives C memory location from R array indices}
\end{Highlighting}
\end{Shaded}

For symmetric matrices and arrays shape can be a pointer to a single
number.

\hypertarget{subroutines}{%
\section{Subroutines}\label{subroutines}}

In array.c there are some auxiliary subroutines for sorting integers
(using the system qsort()) and for binomial coefficients. The main
routines implement encode and decode for arrays with element of type
CTYPE in the following classes.

\begin{itemize}
\tightlist
\item
  general vectors (GVENCODE() and GVDECODE())
\item
  general matrices (GMENCODE() and GMDECODE())
\item
  symmetric matrices in compact storage (SMENCODE() and SMDECODE())
\item
  general multidimensional arrays (GAENCODE() and GADECODE())
\item
  super-symmetric multidimensional arrays in compact storage (SAENCODE()
  and SADECODE())
\end{itemize}

Compact storage, as used here, is column-major and increasing (for
details, see De Leeuw (2017)). Thus array element \(x_{i_1\cdots i_r}\)
is stored if and only if \(i_1\leq\cdots\leq i_r\). The symmetric matrix
versions of decode and encode can also be used for triangular matrices
in compact storage. Of course decode and encode for vectors merely
amounts to subtracting or adding one to the argument, so you want not
want the overhead of a function call in that case.

All \(5\times 2=10\) routines have the same arguments, as in the general
prototypes.

\begin{Shaded}
\begin{Highlighting}[]
\NormalTok{XXDECODE(}\DataTypeTok{const} \DataTypeTok{int}\NormalTok{ *indices, }\DataTypeTok{int}\NormalTok{ *location, }\DataTypeTok{const} \DataTypeTok{int}\NormalTok{ *shape, }\DataTypeTok{const} \DataTypeTok{int}\NormalTok{ *rank);}
\NormalTok{XXENCODE(}\DataTypeTok{int}\NormalTok{ *indices, }\DataTypeTok{const} \DataTypeTok{int}\NormalTok{ *location, }\DataTypeTok{const} \DataTypeTok{int}\NormalTok{ *shape, }\DataTypeTok{const} \DataTypeTok{int}\NormalTok{ *rank);}
\end{Highlighting}
\end{Shaded}

Note that all arguments are pointers to integers. Not all subroutines
need both shape and rank, but we always include them to avoid having to
deal with a variable number of arguments. Note that we only map from R
indices to C storage locations and back, not from C storage locations to
R storage locations, or from C indices to R storage locations. Thus the
workflow we have in mind is creating our data in R, passing it to C,
computing in C using the subroutines, and passing the result back to R.
The 10 subroutines are all written so they can be used directly using
the .C() interface in R (i.e.~they pass all arguments by reference,
return void, and have no I/O). The main reason for not using the .Call()
interface is that we want to make the routines usable for people who do
not have R. We may try at some point to change the R code to use the
.C64() interface (Gerber, Mösinger, and Furrer (2018)).

Using the array structure with pointers to corresponding subroutines
will make it easy to extend our results to other types of arrays, such
as triangular and banded matrices, and even to Hilbert, Vandermonde,
Toeplitz, or Hankel matrices.

\hypertarget{testers}{%
\section{Testers}\label{testers}}

I have also included five main programs in C which test the decode and
encode routines. They are named appropriately:

\begin{itemize}
\tightlist
\item
  gvtester();
\item
  gmtester();
\item
  smtester();
\item
  gatester();
\item
  satester().
\end{itemize}

\hypertarget{appendix-code}{%
\section{Appendix: Code}\label{appendix-code}}

\hypertarget{array.h}{%
\subsection{array.h}\label{array.h}}

\begin{Shaded}
\begin{Highlighting}[]
\PreprocessorTok{\#ifndef INDEXING\_H}
\PreprocessorTok{\#define INDEXING\_H}

\PreprocessorTok{\#include }\ImportTok{\textless{}math.h\textgreater{}}
\PreprocessorTok{\#include }\ImportTok{\textless{}stdio.h\textgreater{}}
\PreprocessorTok{\#include }\ImportTok{\textless{}stdlib.h\textgreater{}}

\PreprocessorTok{\#define CTYPE double}

\KeywordTok{typedef} \DataTypeTok{void}\NormalTok{ (*ENCODE)(}
    \DataTypeTok{const} \DataTypeTok{int}\NormalTok{ *, }\DataTypeTok{int}\NormalTok{ *, }\DataTypeTok{const} \DataTypeTok{int}\NormalTok{ *,}
    \DataTypeTok{const} \DataTypeTok{int}\NormalTok{ *); }\CommentTok{// encode gives R array indices from C memory location}
\KeywordTok{typedef} \DataTypeTok{void}\NormalTok{ (*DECODE)(}
    \DataTypeTok{const} \DataTypeTok{int}\NormalTok{ *, }\DataTypeTok{int}\NormalTok{ *, }\DataTypeTok{const} \DataTypeTok{int}\NormalTok{ *,}
    \DataTypeTok{const} \DataTypeTok{int}\NormalTok{ *); }\CommentTok{// decode gives C memory location from R array indices}

\KeywordTok{typedef} \KeywordTok{struct}\NormalTok{ \{}
\NormalTok{  CTYPE *content;}
  \DataTypeTok{int}\NormalTok{ *length;}
  \DataTypeTok{int}\NormalTok{ *rank;}
  \DataTypeTok{int}\NormalTok{ *shape;}
\NormalTok{  ENCODE encode;}
\NormalTok{  DECODE decode;}
\NormalTok{\} ARRAY;}

\CommentTok{/* protoypes */}

\KeywordTok{extern} \DataTypeTok{int}\NormalTok{ binCoef(}\DataTypeTok{const} \DataTypeTok{int}\NormalTok{, }\DataTypeTok{const} \DataTypeTok{int}\NormalTok{);}

\KeywordTok{extern} \DataTypeTok{void}\NormalTok{ isort(}\DataTypeTok{int}\NormalTok{ *, }\DataTypeTok{const} \DataTypeTok{int}\NormalTok{ *);}

\KeywordTok{extern} \DataTypeTok{int}\NormalTok{ icmp(}\DataTypeTok{const} \DataTypeTok{void}\NormalTok{ *, }\DataTypeTok{const} \DataTypeTok{void}\NormalTok{ *);}

\KeywordTok{extern} \KeywordTok{inline} \DataTypeTok{int}\NormalTok{ IMIN(}\DataTypeTok{const} \DataTypeTok{int}\NormalTok{, }\DataTypeTok{const} \DataTypeTok{int}\NormalTok{);}

\KeywordTok{extern} \DataTypeTok{void}\NormalTok{ GVENCODE(}\DataTypeTok{const} \DataTypeTok{int}\NormalTok{ *, }\DataTypeTok{int}\NormalTok{ *, }\DataTypeTok{const} \DataTypeTok{int}\NormalTok{ *, }\DataTypeTok{const} \DataTypeTok{int}\NormalTok{ *);}

\KeywordTok{extern} \DataTypeTok{void}\NormalTok{ GVDECODE(}\DataTypeTok{const} \DataTypeTok{int}\NormalTok{ *, }\DataTypeTok{int}\NormalTok{ *, }\DataTypeTok{const} \DataTypeTok{int}\NormalTok{ *, }\DataTypeTok{const} \DataTypeTok{int}\NormalTok{ *);}

\KeywordTok{extern} \DataTypeTok{void}\NormalTok{ GMDECODE(}\DataTypeTok{const} \DataTypeTok{int}\NormalTok{ *, }\DataTypeTok{int}\NormalTok{ *, }\DataTypeTok{const} \DataTypeTok{int}\NormalTok{ *, }\DataTypeTok{const} \DataTypeTok{int}\NormalTok{ *);}

\KeywordTok{extern} \DataTypeTok{void}\NormalTok{ GMENCODE(}\DataTypeTok{const} \DataTypeTok{int}\NormalTok{ *, }\DataTypeTok{int}\NormalTok{ *, }\DataTypeTok{const} \DataTypeTok{int}\NormalTok{ *, }\DataTypeTok{const} \DataTypeTok{int}\NormalTok{ *);}

\KeywordTok{extern} \DataTypeTok{void}\NormalTok{ GADECODE(}\DataTypeTok{const} \DataTypeTok{int}\NormalTok{ *, }\DataTypeTok{int}\NormalTok{ *, }\DataTypeTok{const} \DataTypeTok{int}\NormalTok{ *, }\DataTypeTok{const} \DataTypeTok{int}\NormalTok{ *);}

\KeywordTok{extern} \DataTypeTok{void}\NormalTok{ GAENCODE(}\DataTypeTok{const} \DataTypeTok{int}\NormalTok{ *, }\DataTypeTok{int}\NormalTok{ *, }\DataTypeTok{const} \DataTypeTok{int}\NormalTok{ *, }\DataTypeTok{const} \DataTypeTok{int}\NormalTok{ *);}

\KeywordTok{extern} \DataTypeTok{void}\NormalTok{ SMDECODE(}\DataTypeTok{const} \DataTypeTok{int}\NormalTok{ *, }\DataTypeTok{int}\NormalTok{ *, }\DataTypeTok{const} \DataTypeTok{int}\NormalTok{ *, }\DataTypeTok{const} \DataTypeTok{int}\NormalTok{ *);}

\KeywordTok{extern} \DataTypeTok{void}\NormalTok{ SMENCODE(}\DataTypeTok{const} \DataTypeTok{int}\NormalTok{ *, }\DataTypeTok{int}\NormalTok{ *, }\DataTypeTok{const} \DataTypeTok{int}\NormalTok{ *, }\DataTypeTok{const} \DataTypeTok{int}\NormalTok{ *);}

\KeywordTok{extern} \DataTypeTok{void}\NormalTok{ SADECODE(}\DataTypeTok{const} \DataTypeTok{int}\NormalTok{ *, }\DataTypeTok{int}\NormalTok{ *, }\DataTypeTok{const} \DataTypeTok{int}\NormalTok{ *, }\DataTypeTok{const} \DataTypeTok{int}\NormalTok{ *);}

\KeywordTok{extern} \DataTypeTok{void}\NormalTok{ SAENCODE(}\DataTypeTok{const} \DataTypeTok{int}\NormalTok{ *, }\DataTypeTok{int}\NormalTok{ *, }\DataTypeTok{const} \DataTypeTok{int}\NormalTok{ *, }\DataTypeTok{const} \DataTypeTok{int}\NormalTok{ *);}

\PreprocessorTok{\#endif }\CommentTok{/* INDEXING\_H */}
\end{Highlighting}
\end{Shaded}

\hypertarget{array.c}{%
\subsection{array.c}\label{array.c}}

\begin{Shaded}
\begin{Highlighting}[]
\PreprocessorTok{\#include }\ImportTok{"array.h"}

\DataTypeTok{int}\NormalTok{ binCoef(}\DataTypeTok{const} \DataTypeTok{int}\NormalTok{ n, }\DataTypeTok{const} \DataTypeTok{int}\NormalTok{ m) \{}
  \DataTypeTok{int}\NormalTok{ result;}
  \DataTypeTok{int}\NormalTok{ *work = (}\DataTypeTok{int}\NormalTok{ *)calloc((}\DataTypeTok{size\_t}\NormalTok{)(m + }\DecValTok{1}\NormalTok{), }\KeywordTok{sizeof}\NormalTok{(}\DataTypeTok{int}\NormalTok{));}
\NormalTok{  work[}\DecValTok{0}\NormalTok{] = }\DecValTok{1}\NormalTok{;}
  \ControlFlowTok{for}\NormalTok{ (}\DataTypeTok{int}\NormalTok{ i = }\DecValTok{1}\NormalTok{; i \textless{}= n; i++) \{}
    \ControlFlowTok{for}\NormalTok{ (}\DataTypeTok{int}\NormalTok{ j = (i \textgreater{} m) ? m : i; j \textgreater{} }\DecValTok{0}\NormalTok{; j{-}{-}) \{}
\NormalTok{      work[j] += work[j {-} }\DecValTok{1}\NormalTok{];}
\NormalTok{    \}}
\NormalTok{  \}}
\NormalTok{  result = work[m];}
\NormalTok{  free(work);}
  \ControlFlowTok{return}\NormalTok{ (result);}
\NormalTok{\}}

\DataTypeTok{void}\NormalTok{ isort(}\DataTypeTok{int}\NormalTok{ *cell, }\DataTypeTok{const} \DataTypeTok{int}\NormalTok{ *n) \{}
\NormalTok{  (}\DataTypeTok{void}\NormalTok{)qsort(cell, (}\DataTypeTok{size\_t}\NormalTok{)*n, }\KeywordTok{sizeof}\NormalTok{(}\DataTypeTok{int}\NormalTok{), icmp);}
  \ControlFlowTok{return}\NormalTok{;}
\NormalTok{\}}

\DataTypeTok{int}\NormalTok{ icmp(}\DataTypeTok{const} \DataTypeTok{void}\NormalTok{ *x, }\DataTypeTok{const} \DataTypeTok{void}\NormalTok{ *y) \{}
  \DataTypeTok{const} \DataTypeTok{int}\NormalTok{ *ix = (}\DataTypeTok{const} \DataTypeTok{int}\NormalTok{ *)x;}
  \DataTypeTok{const} \DataTypeTok{int}\NormalTok{ *iy = (}\DataTypeTok{const} \DataTypeTok{int}\NormalTok{ *)y;}
  \ControlFlowTok{return}\NormalTok{ (*ix {-} *iy);}
\NormalTok{\}}

\DataTypeTok{void}\NormalTok{ GVENCODE(}\DataTypeTok{const} \DataTypeTok{int}\NormalTok{ *location, }\DataTypeTok{int}\NormalTok{ *indices, }\DataTypeTok{const} \DataTypeTok{int}\NormalTok{ *shape,}
              \DataTypeTok{const} \DataTypeTok{int}\NormalTok{ *rank) \{}
\NormalTok{  *indices = *location + }\DecValTok{1}\NormalTok{;}
  \ControlFlowTok{return}\NormalTok{;}
\NormalTok{\}}

\DataTypeTok{void}\NormalTok{ GVDECODE(}\DataTypeTok{const} \DataTypeTok{int}\NormalTok{ *indices, }\DataTypeTok{int}\NormalTok{ *location, }\DataTypeTok{const} \DataTypeTok{int}\NormalTok{ *shape,}
              \DataTypeTok{const} \DataTypeTok{int}\NormalTok{ *rank) \{}
\NormalTok{  *location = *indices {-} }\DecValTok{1}\NormalTok{;}
  \ControlFlowTok{return}\NormalTok{;}
\NormalTok{\}}

\DataTypeTok{void}\NormalTok{ GMENCODE(}\DataTypeTok{const} \DataTypeTok{int}\NormalTok{ *location, }\DataTypeTok{int}\NormalTok{ *indices, }\DataTypeTok{const} \DataTypeTok{int}\NormalTok{ *shape,}
              \DataTypeTok{const} \DataTypeTok{int}\NormalTok{ *rank) \{}
  \DataTypeTok{int}\NormalTok{ nrow = *(shape + }\DecValTok{0}\NormalTok{), icol, irow, iloc = *location;}
\NormalTok{  irow = (iloc \% nrow) + }\DecValTok{1}\NormalTok{;}
\NormalTok{  icol = (iloc / nrow) + }\DecValTok{1}\NormalTok{;}
\NormalTok{  *(indices + }\DecValTok{0}\NormalTok{) = irow;}
\NormalTok{  *(indices + }\DecValTok{1}\NormalTok{) = icol;}
  \ControlFlowTok{return}\NormalTok{;}
\NormalTok{\}}

\DataTypeTok{void}\NormalTok{ GMDECODE(}\DataTypeTok{const} \DataTypeTok{int}\NormalTok{ *indices, }\DataTypeTok{int}\NormalTok{ *location, }\DataTypeTok{const} \DataTypeTok{int}\NormalTok{ *shape,}
              \DataTypeTok{const} \DataTypeTok{int}\NormalTok{ *rank) \{}
  \DataTypeTok{int}\NormalTok{ nrow = *(shape + }\DecValTok{0}\NormalTok{), irow = *(indices + }\DecValTok{0}\NormalTok{), icol = *(indices + }\DecValTok{1}\NormalTok{),}
\NormalTok{      iloc = (irow {-} }\DecValTok{1}\NormalTok{) + (icol {-} }\DecValTok{1}\NormalTok{) * nrow;}
\NormalTok{  *location = iloc;}
  \ControlFlowTok{return}\NormalTok{;}
\NormalTok{\}}

\DataTypeTok{void}\NormalTok{ SMDECODE(}\DataTypeTok{const} \DataTypeTok{int}\NormalTok{ *indices, }\DataTypeTok{int}\NormalTok{ *location, }\DataTypeTok{const} \DataTypeTok{int}\NormalTok{ *shape,}
              \DataTypeTok{const} \DataTypeTok{int}\NormalTok{ *rank) \{}
  \DataTypeTok{int}\NormalTok{ nrow = *(shape + }\DecValTok{0}\NormalTok{), irow = *(indices + }\DecValTok{0}\NormalTok{), icol = *(indices + }\DecValTok{1}\NormalTok{),}
\NormalTok{      iloc = }\DecValTok{0}\NormalTok{;}
  \ControlFlowTok{if}\NormalTok{ (irow \textless{} icol) \{}
\NormalTok{    iloc = irow;}
\NormalTok{    irow = icol;}
\NormalTok{    icol = iloc;}
\NormalTok{  \}}
  \DataTypeTok{int}\NormalTok{ *maxcol = (}\DataTypeTok{int}\NormalTok{ *)calloc((}\DataTypeTok{size\_t}\NormalTok{)nrow + }\DecValTok{1}\NormalTok{, }\KeywordTok{sizeof}\NormalTok{(}\DataTypeTok{int}\NormalTok{));}
  \ControlFlowTok{for}\NormalTok{ (}\DataTypeTok{int}\NormalTok{ i = }\DecValTok{0}\NormalTok{; i \textless{}= nrow; i++) \{}
\NormalTok{    *(maxcol + i) = i * nrow {-} i * (i {-} }\DecValTok{1}\NormalTok{) / }\DecValTok{2}\NormalTok{;}
\NormalTok{  \}}
\NormalTok{  iloc = *(maxcol + (icol {-} }\DecValTok{1}\NormalTok{)) + (irow {-} icol + }\DecValTok{1}\NormalTok{);}
\NormalTok{  *location = iloc {-} }\DecValTok{1}\NormalTok{;}
\NormalTok{  free(maxcol);}
  \ControlFlowTok{return}\NormalTok{;}
\NormalTok{\}}

\DataTypeTok{void}\NormalTok{ SMENCODE(}\DataTypeTok{const} \DataTypeTok{int}\NormalTok{ *location, }\DataTypeTok{int}\NormalTok{ *indices, }\DataTypeTok{const} \DataTypeTok{int}\NormalTok{ *shape,}
              \DataTypeTok{const} \DataTypeTok{int}\NormalTok{ *rank) \{}
  \DataTypeTok{int}\NormalTok{ nrow = *(shape + }\DecValTok{0}\NormalTok{), irow, icol, iloc = *location + }\DecValTok{1}\NormalTok{;}
  \DataTypeTok{int}\NormalTok{ *maxcol = (}\DataTypeTok{int}\NormalTok{ *)calloc((}\DataTypeTok{size\_t}\NormalTok{)nrow + }\DecValTok{1}\NormalTok{, }\KeywordTok{sizeof}\NormalTok{(}\DataTypeTok{int}\NormalTok{));}
  \ControlFlowTok{for}\NormalTok{ (}\DataTypeTok{int}\NormalTok{ i = }\DecValTok{0}\NormalTok{; i \textless{}= nrow; i++) \{}
\NormalTok{    *(maxcol + i) = i * nrow {-} i * (i {-} }\DecValTok{1}\NormalTok{) / }\DecValTok{2}\NormalTok{;}
\NormalTok{  \}}
  \ControlFlowTok{for}\NormalTok{ (}\DataTypeTok{int}\NormalTok{ i = }\DecValTok{1}\NormalTok{; i \textless{}= nrow; i++) \{}
    \ControlFlowTok{if}\NormalTok{ (iloc \textless{}= *(maxcol + i)) \{}
\NormalTok{      icol = i;}
\NormalTok{      irow = iloc {-} *(maxcol + (i {-} }\DecValTok{1}\NormalTok{)) + (i {-} }\DecValTok{1}\NormalTok{);}
      \ControlFlowTok{break}\NormalTok{;}
\NormalTok{    \}}
\NormalTok{  \}}
\NormalTok{  *(indices + }\DecValTok{0}\NormalTok{) = irow;}
\NormalTok{  *(indices + }\DecValTok{1}\NormalTok{) = icol;}
\NormalTok{  free(maxcol);}
  \ControlFlowTok{return}\NormalTok{;}
\NormalTok{\}}

\DataTypeTok{void}\NormalTok{ GADECODE(}\DataTypeTok{const} \DataTypeTok{int}\NormalTok{ *indices, }\DataTypeTok{int}\NormalTok{ *location, }\DataTypeTok{const} \DataTypeTok{int}\NormalTok{ *shape,}
              \DataTypeTok{const} \DataTypeTok{int}\NormalTok{ *rank) \{}
  \DataTypeTok{int}\NormalTok{ aux = }\DecValTok{1}\NormalTok{, res = }\DecValTok{1}\NormalTok{, r = *rank;}
  \ControlFlowTok{for}\NormalTok{ (}\DataTypeTok{int}\NormalTok{ i = }\DecValTok{0}\NormalTok{; i \textless{} r; i++) \{}
\NormalTok{    res += (*(indices + i) {-} }\DecValTok{1}\NormalTok{) * aux;}
\NormalTok{    aux *= *(shape + i);}
\NormalTok{  \}}
\NormalTok{  *location = res {-} }\DecValTok{1}\NormalTok{;}
  \ControlFlowTok{return}\NormalTok{;}
\NormalTok{\}}

\DataTypeTok{void}\NormalTok{ GAENCODE(}\DataTypeTok{const} \DataTypeTok{int}\NormalTok{ *location, }\DataTypeTok{int}\NormalTok{ *indices, }\DataTypeTok{const} \DataTypeTok{int}\NormalTok{ *shape,}
              \DataTypeTok{const} \DataTypeTok{int}\NormalTok{ *rank) \{}
  \DataTypeTok{int}\NormalTok{ r = *rank, iloc = *location + }\DecValTok{1}\NormalTok{, aux = }\DecValTok{1}\NormalTok{, k = }\DecValTok{1}\NormalTok{;}
  \ControlFlowTok{for}\NormalTok{ (}\DataTypeTok{int}\NormalTok{ j = }\DecValTok{1}\NormalTok{; j \textless{} r; j++) \{}
\NormalTok{    aux *= *(shape + j {-} }\DecValTok{1}\NormalTok{);}
\NormalTok{  \}}
  \ControlFlowTok{for}\NormalTok{ (}\DataTypeTok{int}\NormalTok{ j = r {-} }\DecValTok{1}\NormalTok{; j \textgreater{} }\DecValTok{0}\NormalTok{; j{-}{-}) \{}
\NormalTok{    k = (iloc {-} }\DecValTok{1}\NormalTok{) / aux;}
\NormalTok{    *(indices + j) = k + }\DecValTok{1}\NormalTok{;}
\NormalTok{    iloc {-}= k * aux;}
\NormalTok{    aux /= *(shape + j {-} }\DecValTok{1}\NormalTok{);}
\NormalTok{  \}}
\NormalTok{  *(indices + }\DecValTok{0}\NormalTok{) = iloc;}
  \ControlFlowTok{return}\NormalTok{;}
\NormalTok{\}}

\DataTypeTok{void}\NormalTok{ SADECODE(}\DataTypeTok{const} \DataTypeTok{int}\NormalTok{ *indices, }\DataTypeTok{int}\NormalTok{ *location, }\DataTypeTok{const} \DataTypeTok{int}\NormalTok{ *shape,}
              \DataTypeTok{const} \DataTypeTok{int}\NormalTok{ *rank) \{}
  \DataTypeTok{int}\NormalTok{ iloc = }\DecValTok{1}\NormalTok{, k = }\DecValTok{1}\NormalTok{, r = *rank;}
  \DataTypeTok{int}\NormalTok{ *ind = (}\DataTypeTok{int}\NormalTok{ *)calloc((}\DataTypeTok{size\_t}\NormalTok{)r, }\KeywordTok{sizeof}\NormalTok{(}\DataTypeTok{int}\NormalTok{));}
  \ControlFlowTok{for}\NormalTok{ (}\DataTypeTok{int}\NormalTok{ i = }\DecValTok{0}\NormalTok{; i \textless{} r; i++) \{}
\NormalTok{    *(ind + i) = *(indices + i);}
\NormalTok{  \}}
\NormalTok{  (}\DataTypeTok{void}\NormalTok{)isort(ind, rank);}
  \ControlFlowTok{for}\NormalTok{ (}\DataTypeTok{int}\NormalTok{ i = }\DecValTok{1}\NormalTok{; i \textless{}= r; i++) \{}
\NormalTok{    k = i + (*(ind + i {-} }\DecValTok{1}\NormalTok{) {-} }\DecValTok{1}\NormalTok{) {-} }\DecValTok{1}\NormalTok{;}
\NormalTok{    iloc += binCoef(k, i);}
\NormalTok{  \}}
\NormalTok{  *location = iloc {-} }\DecValTok{1}\NormalTok{;}
\NormalTok{  (}\DataTypeTok{void}\NormalTok{)free(ind);}
  \ControlFlowTok{return}\NormalTok{;}
\NormalTok{\}}

\DataTypeTok{void}\NormalTok{ SAENCODE(}\DataTypeTok{const} \DataTypeTok{int}\NormalTok{ *location, }\DataTypeTok{int}\NormalTok{ *indices, }\DataTypeTok{const} \DataTypeTok{int}\NormalTok{ *shape,}
              \DataTypeTok{const} \DataTypeTok{int}\NormalTok{ *rank) \{}
  \DataTypeTok{int}\NormalTok{ n = *shape, r = *rank, iloc = *location + }\DecValTok{1}\NormalTok{, v = iloc {-} }\DecValTok{1}\NormalTok{;}
  \ControlFlowTok{for}\NormalTok{ (}\DataTypeTok{int}\NormalTok{ k = r; k \textgreater{}= }\DecValTok{1}\NormalTok{; k{-}{-}) \{}
    \ControlFlowTok{for}\NormalTok{ (}\DataTypeTok{int}\NormalTok{ j = }\DecValTok{0}\NormalTok{; j \textless{} n; j++) \{}
      \DataTypeTok{int}\NormalTok{ sj = binCoef(k + j {-} }\DecValTok{1}\NormalTok{, k), sk = binCoef(k + j, k);}
      \ControlFlowTok{if}\NormalTok{ (v \textless{} sk) \{}
\NormalTok{        indices[k {-} }\DecValTok{1}\NormalTok{] = j + }\DecValTok{1}\NormalTok{;}
\NormalTok{        v {-}= sj;}
        \ControlFlowTok{break}\NormalTok{;}
\NormalTok{      \}}
\NormalTok{    \}}
\NormalTok{  \}}
  \ControlFlowTok{return}\NormalTok{;}
\NormalTok{\}}
\end{Highlighting}
\end{Shaded}

\hypertarget{gvtester.c}{%
\subsection{gvtester.c}\label{gvtester.c}}

\begin{Shaded}
\begin{Highlighting}[]
\PreprocessorTok{\#include }\ImportTok{"array.h"}

\CommentTok{/*}
\CommentTok{  a vector of shape (10)}
\CommentTok{*/}

\DataTypeTok{int}\NormalTok{ main(}\DataTypeTok{void}\NormalTok{) \{}
\NormalTok{  ARRAY a;}
  \DataTypeTok{int}\NormalTok{ length = }\DecValTok{10}\NormalTok{, rank = }\DecValTok{1}\NormalTok{, j = }\DecValTok{0}\NormalTok{;}
  \DataTypeTok{int}\NormalTok{ *shape = (}\DataTypeTok{int}\NormalTok{ *)calloc((}\DataTypeTok{size\_t}\NormalTok{)rank, }\KeywordTok{sizeof}\NormalTok{(}\DataTypeTok{int}\NormalTok{));}
\NormalTok{  *shape = }\DecValTok{10}\NormalTok{;}
\NormalTok{  a.content = (}\DataTypeTok{double}\NormalTok{ *)calloc((}\DataTypeTok{size\_t}\NormalTok{)length, }\KeywordTok{sizeof}\NormalTok{(}\DataTypeTok{double}\NormalTok{));}
\NormalTok{  a.length = \&length;}
\NormalTok{  a.shape = shape;}
\NormalTok{  a.rank = \&rank;}
\NormalTok{  a.decode = GVDECODE;}
\NormalTok{  a.encode = GVENCODE;}
\NormalTok{  printf(}\StringTok{"========}\SpecialCharTok{\textbackslash{}n}\StringTok{"}\NormalTok{);}
\NormalTok{  printf(}\StringTok{"GVDECODE}\SpecialCharTok{\textbackslash{}n}\StringTok{"}\NormalTok{);}
\NormalTok{  printf(}\StringTok{"========}\SpecialCharTok{\textbackslash{}n}\StringTok{"}\NormalTok{);}
  \ControlFlowTok{for}\NormalTok{ (}\DataTypeTok{int}\NormalTok{ i = }\DecValTok{1}\NormalTok{; i \textless{}= length; i++) \{}
\NormalTok{    (}\DataTypeTok{void}\NormalTok{)a.decode(\&i, \&j, a.shape, a.rank);}
\NormalTok{    printf(}\StringTok{"index \%4d location \%4d}\SpecialCharTok{\textbackslash{}n}\StringTok{"}\NormalTok{, i, j);}
\NormalTok{  \}}
\NormalTok{  printf(}\StringTok{"========}\SpecialCharTok{\textbackslash{}n}\StringTok{"}\NormalTok{);}
\NormalTok{  printf(}\StringTok{"GVENCODE}\SpecialCharTok{\textbackslash{}n}\StringTok{"}\NormalTok{);}
\NormalTok{  printf(}\StringTok{"========}\SpecialCharTok{\textbackslash{}n}\StringTok{"}\NormalTok{);}
  \ControlFlowTok{for}\NormalTok{ (}\DataTypeTok{int}\NormalTok{ i = }\DecValTok{0}\NormalTok{; i \textless{} length; i++) \{}
\NormalTok{    (}\DataTypeTok{void}\NormalTok{)a.encode(\&i, \&j, a.shape, a.rank);}
\NormalTok{    printf(}\StringTok{"location \%4d index \%4d}\SpecialCharTok{\textbackslash{}n}\StringTok{"}\NormalTok{, i, j);}
\NormalTok{  \}}
\NormalTok{  (}\DataTypeTok{void}\NormalTok{)free(shape);}
\NormalTok{  (}\DataTypeTok{void}\NormalTok{)free(a.content);}
  \ControlFlowTok{return}\NormalTok{ EXIT\_SUCCESS;}
\NormalTok{\}}
\end{Highlighting}
\end{Shaded}

\hypertarget{gmtester.c}{%
\subsection{gmtester.c}\label{gmtester.c}}

\begin{Shaded}
\begin{Highlighting}[]
\PreprocessorTok{\#include }\ImportTok{"array.h"}

\CommentTok{/*}
\CommentTok{  a general matrix of shape (4,3)}
\CommentTok{*/}

\DataTypeTok{int}\NormalTok{ main(}\DataTypeTok{void}\NormalTok{) \{}
\NormalTok{  ARRAY a;}
  \DataTypeTok{int}\NormalTok{ length = }\DecValTok{12}\NormalTok{, rank = }\DecValTok{2}\NormalTok{, k = }\DecValTok{0}\NormalTok{;}
  \DataTypeTok{int}\NormalTok{ *shape = (}\DataTypeTok{int}\NormalTok{ *)calloc((}\DataTypeTok{size\_t}\NormalTok{)rank, }\KeywordTok{sizeof}\NormalTok{(}\DataTypeTok{int}\NormalTok{));}
  \DataTypeTok{int}\NormalTok{ *indices = (}\DataTypeTok{int}\NormalTok{ *)calloc((}\DataTypeTok{size\_t}\NormalTok{)rank, }\KeywordTok{sizeof}\NormalTok{(}\DataTypeTok{int}\NormalTok{));}
\NormalTok{  *(shape + }\DecValTok{0}\NormalTok{) = }\DecValTok{4}\NormalTok{;}
\NormalTok{  *(shape + }\DecValTok{1}\NormalTok{) = }\DecValTok{3}\NormalTok{;}
\NormalTok{  a.content = (}\DataTypeTok{double}\NormalTok{ *)calloc((}\DataTypeTok{size\_t}\NormalTok{)length, }\KeywordTok{sizeof}\NormalTok{(}\DataTypeTok{double}\NormalTok{));}
\NormalTok{  a.length = \&length;}
\NormalTok{  a.shape = shape;}
\NormalTok{  a.rank = \&rank;}
\NormalTok{  a.decode = GMDECODE;}
\NormalTok{  a.encode = GMENCODE;}
\NormalTok{  printf(}\StringTok{"========}\SpecialCharTok{\textbackslash{}n}\StringTok{"}\NormalTok{);}
\NormalTok{  printf(}\StringTok{"GMDECODE}\SpecialCharTok{\textbackslash{}n}\StringTok{"}\NormalTok{);}
\NormalTok{  printf(}\StringTok{"========}\SpecialCharTok{\textbackslash{}n}\StringTok{"}\NormalTok{);}
  \ControlFlowTok{for}\NormalTok{ (}\DataTypeTok{int}\NormalTok{ j = }\DecValTok{1}\NormalTok{; j \textless{}= *(a.shape + }\DecValTok{1}\NormalTok{); j++) \{}
\NormalTok{    *(indices + }\DecValTok{1}\NormalTok{) = j;}
    \ControlFlowTok{for}\NormalTok{ (}\DataTypeTok{int}\NormalTok{ i = }\DecValTok{1}\NormalTok{; i \textless{}= *(a.shape + }\DecValTok{0}\NormalTok{); i++) \{}
\NormalTok{      *(indices + }\DecValTok{0}\NormalTok{) = i;}
\NormalTok{      (}\DataTypeTok{void}\NormalTok{)a.decode(indices, \&k, a.shape, a.rank);}
\NormalTok{      printf(}\StringTok{"indices \%4d \%4d location \%4d}\SpecialCharTok{\textbackslash{}n}\StringTok{"}\NormalTok{, i, j, k);}
\NormalTok{    \}}
\NormalTok{  \}}
\NormalTok{  printf(}\StringTok{"========}\SpecialCharTok{\textbackslash{}n}\StringTok{"}\NormalTok{);}
\NormalTok{  printf(}\StringTok{"GMENCODE}\SpecialCharTok{\textbackslash{}n}\StringTok{"}\NormalTok{);}
\NormalTok{  printf(}\StringTok{"========}\SpecialCharTok{\textbackslash{}n}\StringTok{"}\NormalTok{);}
  \ControlFlowTok{for}\NormalTok{ (}\DataTypeTok{int}\NormalTok{ i = }\DecValTok{0}\NormalTok{; i \textless{} length; i++) \{}
\NormalTok{    (}\DataTypeTok{void}\NormalTok{)a.encode(\&i, indices, a.shape, a.rank);}
\NormalTok{    printf(}\StringTok{"location \%4d indices \%4d \%4d}\SpecialCharTok{\textbackslash{}n}\StringTok{"}\NormalTok{, i, *(indices + }\DecValTok{0}\NormalTok{), *(indices + }\DecValTok{1}\NormalTok{));}
\NormalTok{  \}}
\NormalTok{  (}\DataTypeTok{void}\NormalTok{)free(shape);}
\NormalTok{  (}\DataTypeTok{void}\NormalTok{)free(a.content);}
\NormalTok{  (}\DataTypeTok{void}\NormalTok{)free(indices);}
  \ControlFlowTok{return}\NormalTok{ EXIT\_SUCCESS;}
\NormalTok{\}}
\end{Highlighting}
\end{Shaded}

\hypertarget{smtester.c}{%
\subsection{smtester.c}\label{smtester.c}}

\begin{Shaded}
\begin{Highlighting}[]
\PreprocessorTok{\#include }\ImportTok{"array.h"}

\CommentTok{/*}
\CommentTok{  a symmetric matrix of shape (4,4)}
\CommentTok{*/}

\DataTypeTok{int}\NormalTok{ main(}\DataTypeTok{void}\NormalTok{) \{}
\NormalTok{  ARRAY a;}
  \DataTypeTok{int}\NormalTok{ length = }\DecValTok{10}\NormalTok{, rank = }\DecValTok{2}\NormalTok{, shape = }\DecValTok{4}\NormalTok{, k = }\DecValTok{0}\NormalTok{;}
  \DataTypeTok{int}\NormalTok{ *indices = (}\DataTypeTok{int}\NormalTok{ *)calloc((}\DataTypeTok{size\_t}\NormalTok{)rank, }\KeywordTok{sizeof}\NormalTok{(}\DataTypeTok{int}\NormalTok{));}
\NormalTok{  a.content = (CTYPE *)calloc((}\DataTypeTok{size\_t}\NormalTok{)length, }\KeywordTok{sizeof}\NormalTok{(CTYPE));}
\NormalTok{  a.length = \&length;}
\NormalTok{  a.shape = \&shape;}
\NormalTok{  a.rank = \&rank;}
\NormalTok{  a.decode = SMDECODE;}
\NormalTok{  a.encode = SMENCODE;}
\NormalTok{  printf(}\StringTok{"========}\SpecialCharTok{\textbackslash{}n}\StringTok{"}\NormalTok{);}
\NormalTok{  printf(}\StringTok{"SMDECODE}\SpecialCharTok{\textbackslash{}n}\StringTok{"}\NormalTok{);}
\NormalTok{  printf(}\StringTok{"========}\SpecialCharTok{\textbackslash{}n}\StringTok{"}\NormalTok{);}
  \ControlFlowTok{for}\NormalTok{ (}\DataTypeTok{int}\NormalTok{ j = }\DecValTok{1}\NormalTok{; j \textless{}= shape; j++) \{}
\NormalTok{    *(indices + }\DecValTok{1}\NormalTok{) = j;}
    \ControlFlowTok{for}\NormalTok{ (}\DataTypeTok{int}\NormalTok{ i = }\DecValTok{1}\NormalTok{; i \textless{}= shape; i++) \{}
\NormalTok{      *(indices + }\DecValTok{0}\NormalTok{) = i;}
\NormalTok{      (}\DataTypeTok{void}\NormalTok{)a.decode(indices, \&k, a.shape, a.rank);}
\NormalTok{      printf(}\StringTok{"indices \%4d \%4d location \%4d}\SpecialCharTok{\textbackslash{}n}\StringTok{"}\NormalTok{, i, j, k);}
\NormalTok{    \}}
\NormalTok{  \}}
\NormalTok{  printf(}\StringTok{"========}\SpecialCharTok{\textbackslash{}n}\StringTok{"}\NormalTok{);}
\NormalTok{  printf(}\StringTok{"SMENCODE}\SpecialCharTok{\textbackslash{}n}\StringTok{"}\NormalTok{);}
\NormalTok{  printf(}\StringTok{"========}\SpecialCharTok{\textbackslash{}n}\StringTok{"}\NormalTok{);}
  \ControlFlowTok{for}\NormalTok{ (}\DataTypeTok{int}\NormalTok{ i = }\DecValTok{0}\NormalTok{; i \textless{} length; i++) \{}
\NormalTok{    (}\DataTypeTok{void}\NormalTok{)a.encode(\&i, indices, a.shape, a.rank);}
\NormalTok{    printf(}\StringTok{"location \%4d indices \%4d \%4d}\SpecialCharTok{\textbackslash{}n}\StringTok{"}\NormalTok{, i, *(indices + }\DecValTok{0}\NormalTok{), *(indices + }\DecValTok{1}\NormalTok{));}
\NormalTok{  \}}
\NormalTok{  (}\DataTypeTok{void}\NormalTok{)free(a.content);}
\NormalTok{  (}\DataTypeTok{void}\NormalTok{)free(indices);}
  \ControlFlowTok{return}\NormalTok{ EXIT\_SUCCESS;}
\NormalTok{\}}
\end{Highlighting}
\end{Shaded}

\hypertarget{gatester.c}{%
\subsection{gatester.c}\label{gatester.c}}

\begin{Shaded}
\begin{Highlighting}[]
\PreprocessorTok{\#include }\ImportTok{"array.h"}

\CommentTok{/*}
\CommentTok{  a general array of shape (2,4,3)}
\CommentTok{*/}

\DataTypeTok{int}\NormalTok{ main(}\DataTypeTok{void}\NormalTok{) \{}
\NormalTok{  ARRAY a;}
  \DataTypeTok{int}\NormalTok{ length = }\DecValTok{24}\NormalTok{, rank = }\DecValTok{3}\NormalTok{, l = }\DecValTok{0}\NormalTok{;}
  \DataTypeTok{int}\NormalTok{ *shape = (}\DataTypeTok{int}\NormalTok{ *)calloc((}\DataTypeTok{size\_t}\NormalTok{)rank, }\KeywordTok{sizeof}\NormalTok{(}\DataTypeTok{int}\NormalTok{));}
  \DataTypeTok{int}\NormalTok{ *indices = (}\DataTypeTok{int}\NormalTok{ *)calloc((}\DataTypeTok{size\_t}\NormalTok{)rank, }\KeywordTok{sizeof}\NormalTok{(}\DataTypeTok{int}\NormalTok{));}
\NormalTok{  *(shape + }\DecValTok{0}\NormalTok{) = }\DecValTok{2}\NormalTok{;}
\NormalTok{  *(shape + }\DecValTok{1}\NormalTok{) = }\DecValTok{4}\NormalTok{;}
\NormalTok{  *(shape + }\DecValTok{2}\NormalTok{) = }\DecValTok{3}\NormalTok{;}
\NormalTok{  a.content = (}\DataTypeTok{double}\NormalTok{ *)calloc((}\DataTypeTok{size\_t}\NormalTok{)length, }\KeywordTok{sizeof}\NormalTok{(}\DataTypeTok{double}\NormalTok{));}
\NormalTok{  a.length = \&length;}
\NormalTok{  a.shape = shape;}
\NormalTok{  a.rank = \&rank;}
\NormalTok{  a.decode = GADECODE;}
\NormalTok{  a.encode = GAENCODE;}
\NormalTok{  printf(}\StringTok{"========}\SpecialCharTok{\textbackslash{}n}\StringTok{"}\NormalTok{);}
\NormalTok{  printf(}\StringTok{"GADECODE}\SpecialCharTok{\textbackslash{}n}\StringTok{"}\NormalTok{);}
\NormalTok{  printf(}\StringTok{"========}\SpecialCharTok{\textbackslash{}n}\StringTok{"}\NormalTok{);}
  \ControlFlowTok{for}\NormalTok{ (}\DataTypeTok{int}\NormalTok{ k = }\DecValTok{1}\NormalTok{; k \textless{}= *(a.shape + }\DecValTok{2}\NormalTok{); k++) \{}
\NormalTok{    *(indices + }\DecValTok{2}\NormalTok{) = k;}
    \ControlFlowTok{for}\NormalTok{ (}\DataTypeTok{int}\NormalTok{ j = }\DecValTok{1}\NormalTok{; j \textless{}= *(a.shape + }\DecValTok{1}\NormalTok{); j++) \{}
\NormalTok{      *(indices + }\DecValTok{1}\NormalTok{) = j;}
      \ControlFlowTok{for}\NormalTok{ (}\DataTypeTok{int}\NormalTok{ i = }\DecValTok{1}\NormalTok{; i \textless{}= *(a.shape + }\DecValTok{0}\NormalTok{); i++) \{}
\NormalTok{        *(indices + }\DecValTok{0}\NormalTok{) = i;}
\NormalTok{        (}\DataTypeTok{void}\NormalTok{)a.decode(indices, \&l, a.shape, a.rank);}
\NormalTok{        printf(}\StringTok{"indices \%4d \%4d \%4d location \%4d}\SpecialCharTok{\textbackslash{}n}\StringTok{"}\NormalTok{, i, j, k, l);}
\NormalTok{      \}}
\NormalTok{    \}}
\NormalTok{  \}}
\NormalTok{  printf(}\StringTok{"========}\SpecialCharTok{\textbackslash{}n}\StringTok{"}\NormalTok{);}
\NormalTok{  printf(}\StringTok{"GAENCODE}\SpecialCharTok{\textbackslash{}n}\StringTok{"}\NormalTok{);}
\NormalTok{  printf(}\StringTok{"========}\SpecialCharTok{\textbackslash{}n}\StringTok{"}\NormalTok{);}
  \ControlFlowTok{for}\NormalTok{ (}\DataTypeTok{int}\NormalTok{ i = }\DecValTok{0}\NormalTok{; i \textless{} length; i++) \{}
\NormalTok{    (}\DataTypeTok{void}\NormalTok{)a.encode(\&i, indices, a.shape, a.rank);}
\NormalTok{    printf(}\StringTok{"location \%4d indices \%4d \%4d \%4d}\SpecialCharTok{\textbackslash{}n}\StringTok{"}\NormalTok{, i, *(indices + }\DecValTok{0}\NormalTok{),}
\NormalTok{           *(indices + }\DecValTok{1}\NormalTok{), *(indices + }\DecValTok{2}\NormalTok{));}
\NormalTok{  \}}
\NormalTok{  (}\DataTypeTok{void}\NormalTok{)free(shape);}
\NormalTok{  (}\DataTypeTok{void}\NormalTok{)free(a.content);}
\NormalTok{  (}\DataTypeTok{void}\NormalTok{)free(indices);}
\NormalTok{\}}
\end{Highlighting}
\end{Shaded}

\hypertarget{satester.c}{%
\subsection{satester.c}\label{satester.c}}

\begin{Shaded}
\begin{Highlighting}[]
\PreprocessorTok{\#include }\ImportTok{"array.h"}

\CommentTok{/*}
\CommentTok{  a symmetric array of shape (3,3,3)}
\CommentTok{*/}

\DataTypeTok{int}\NormalTok{ main(}\DataTypeTok{void}\NormalTok{) \{}
\NormalTok{  ARRAY a;}
  \DataTypeTok{int}\NormalTok{ length = }\DecValTok{10}\NormalTok{, rank = }\DecValTok{3}\NormalTok{, shape = }\DecValTok{3}\NormalTok{, l = }\DecValTok{0}\NormalTok{;}
  \DataTypeTok{int}\NormalTok{ *indices = (}\DataTypeTok{int}\NormalTok{ *)calloc((}\DataTypeTok{size\_t}\NormalTok{)rank, }\KeywordTok{sizeof}\NormalTok{(}\DataTypeTok{int}\NormalTok{));}
\NormalTok{  a.content = (}\DataTypeTok{double}\NormalTok{ *)calloc((}\DataTypeTok{size\_t}\NormalTok{)length, }\KeywordTok{sizeof}\NormalTok{(}\DataTypeTok{double}\NormalTok{));}
\NormalTok{  a.length = \&length;}
\NormalTok{  a.shape = \&shape;}
\NormalTok{  a.rank = \&rank;}
\NormalTok{  a.decode = SADECODE;}
\NormalTok{  a.encode = SAENCODE;}
\NormalTok{  printf(}\StringTok{"========}\SpecialCharTok{\textbackslash{}n}\StringTok{"}\NormalTok{);}
\NormalTok{  printf(}\StringTok{"SADECODE}\SpecialCharTok{\textbackslash{}n}\StringTok{"}\NormalTok{);}
\NormalTok{  printf(}\StringTok{"========}\SpecialCharTok{\textbackslash{}n}\StringTok{"}\NormalTok{);}
  \ControlFlowTok{for}\NormalTok{ (}\DataTypeTok{int}\NormalTok{ k = }\DecValTok{1}\NormalTok{; k \textless{}= shape; k++) \{}
\NormalTok{    *(indices + }\DecValTok{2}\NormalTok{) = k;}
    \ControlFlowTok{for}\NormalTok{ (}\DataTypeTok{int}\NormalTok{ j = }\DecValTok{1}\NormalTok{; j \textless{}= shape; j++) \{}
\NormalTok{      *(indices + }\DecValTok{1}\NormalTok{) = j;}
      \ControlFlowTok{for}\NormalTok{ (}\DataTypeTok{int}\NormalTok{ i = }\DecValTok{1}\NormalTok{; i \textless{}= shape; i++) \{}
\NormalTok{        *(indices + }\DecValTok{0}\NormalTok{) = i;}
        \ControlFlowTok{if}\NormalTok{ ((i \textless{}= j) \&\& (j \textless{}= k)) \{}
\NormalTok{          (}\DataTypeTok{void}\NormalTok{)a.decode(indices, \&l, a.shape, a.rank);}
\NormalTok{          printf(}\StringTok{"indices \%4d \%4d \%4d location \%4d}\SpecialCharTok{\textbackslash{}n}\StringTok{"}\NormalTok{, i, j, k, l);}
\NormalTok{        \}}
\NormalTok{      \}}
\NormalTok{    \}}
\NormalTok{  \}}
\NormalTok{  printf(}\StringTok{"========}\SpecialCharTok{\textbackslash{}n}\StringTok{"}\NormalTok{);}
\NormalTok{  printf(}\StringTok{"SAENCODE}\SpecialCharTok{\textbackslash{}n}\StringTok{"}\NormalTok{);}
\NormalTok{  printf(}\StringTok{"========}\SpecialCharTok{\textbackslash{}n}\StringTok{"}\NormalTok{);}
  \ControlFlowTok{for}\NormalTok{ (}\DataTypeTok{int}\NormalTok{ i = }\DecValTok{0}\NormalTok{; i \textless{} length; i++) \{}
\NormalTok{    (}\DataTypeTok{void}\NormalTok{)a.encode(\&i, indices, a.shape, a.rank);}
\NormalTok{    printf(}\StringTok{"location \%4d indices \%4d \%4d \%4d}\SpecialCharTok{\textbackslash{}n}\StringTok{"}\NormalTok{, i, *(indices + }\DecValTok{0}\NormalTok{),}
\NormalTok{           *(indices + }\DecValTok{1}\NormalTok{), *(indices + }\DecValTok{2}\NormalTok{));}
\NormalTok{  \}}
\NormalTok{  (}\DataTypeTok{void}\NormalTok{)free(a.content);}
\NormalTok{  (}\DataTypeTok{void}\NormalTok{)free(indices);}
  \ControlFlowTok{return}\NormalTok{ EXIT\_SUCCESS;}
\NormalTok{\}}
\end{Highlighting}
\end{Shaded}

\hypertarget{references}{%
\section*{References}\label{references}}
\addcontentsline{toc}{section}{References}

\hypertarget{refs}{}
\begin{cslreferences}
\leavevmode\hypertarget{ref-deleeuw_E_17s}{}%
De Leeuw, J. 2017. ``Multidimensional Array Indexing and Storage.''
2017. \url{http://deleeuwpdx.net/pubfolders/indexing/indexing.pdf}.

\leavevmode\hypertarget{ref-gerber_moesinger_furrer_18}{}%
Gerber, F., K. Mösinger, and R. Furrer. 2018. ``dotCall64: An efficient
interface to compiled C/C++ and Fortran code supporting long vectors.''
\emph{SoftwareX} 7: 217--21.
\end{cslreferences}

\end{document}
